\section{Cost Model}
\label{sec:cost_model}

For a run $r$ of strategy $s$ at training size $n$, let
$S_{r,s,n}\subseteq C$ be the selected sample. The usual cost model is the
count of measured configurations:
\begin{equation}
\SizeCost(s,n)=n.
\label{eq:sizecost}
\end{equation}
We compare it with the measured cost:
\begin{equation}
\TrainCost(r,s,n)=\sum_{c\in S_{r,s,n}}\mathit{cost}(c).
\label{eq:traincost}
\end{equation}
Here $\mathit{cost}(c)$ is the measured resource consumed for configuration $c$:
wall-clock measurement time in most of our study, energy in the
performance--energy family, and the uncompressed \texttt{vmlinux} size in the
Linux family. The first two are label-acquisition costs. Image size is instead a
resource budget (the shippable footprint that constrains embedded and
over-the-air deployment), so we rest the acquisition-cost claim on the five
families where cost is acquisition cost and treat Linux as a resource-budget
stress test that isolates strategy-induced selection bias alone: there the
prediction target and the paid resource coincide, leaving no target--cost
asymmetry to exploit (\Cref{subsec:results-rq3}). The model does not require the
cost dimension to equal the prediction target: one may predict output size while
paying time, or runtime while paying energy.

To compare systems with different absolute costs, we normalize by the cost of
measuring the whole available configuration space:
\begin{align}
\FullCost &= \sum_{c\in C}\mathit{cost}(c),\\
\mathit{TrainCostPct} &= \frac{\TrainCost}{\FullCost},&
\mathit{SizePct} &= \frac{n}{|C|}.
\label{eq:pct}
\end{align}
$\FullCost$ is a retrospective normalizer for cross-system reporting only: the
$\Ratio$ and percentage views use it, whereas the strategy rankings and Pareto
fronts compare raw measured $\TrainCost$. A deployable sampler therefore never
needs $\FullCost$ in advance.
The count proxy has good \emph{budget agreement} when the fraction of
configurations measured matches the fraction of measured cost spent. We also
call this \emph{cost-proxy calibration}: the sample-size proxy is calibrated
when a $10\%$ sample spends about $10\%$ of the measured budget. This is not
statistical calibration of a prediction model. It is only a check that the cost
proxy is on the same scale as the measured resource. We summarize it with a
budget ratio:
\begin{equation}
\Ratio(r,s,n)=
\frac{\mathit{TrainCostPct}(r,s,n)}{\mathit{SizePct}(s,n)}.
\label{eq:ratio}
\end{equation}
$\Ratio=1$ means that a $10\%$ sample consumes about $10\%$ of the full
measurement cost. $\Ratio>1$ means the strategy selected costlier-than-average
configurations, whereas $\Ratio<1$ means it selected cheaper-than-average ones. This is
a per-run ratio, so one-number summaries report the median ratio unless stated
otherwise. The practical question is whether a same-size sample takes more or
less of the real budget than the count suggests.

For decision-level comparisons, we use two standard views. First, we rank
strategy--budget combinations, where each combination is one sampling strategy
at one budget level, with a normalized objective
\begin{align}
J^{\text{size}}_{\lambda} &= \widetilde{\MRE}+\lambda\widetilde{\SizeCost},\\
J^{\text{meas}}_{\lambda} &= \widetilde{\MRE}+\lambda\widetilde{\TrainCost},
\end{align}
where the cost weight $\lambda\in\{0.1,0.25,0.5,1,2,4,10\}$ spans
accuracy-heavy to cost-heavy choices. Second, we compare Pareto fronts in the
MRE--cost plane: a combination is efficient if no other combination has both
lower error and lower cost. The main threat we test is whether the recommended
strategy or the efficient set changes when $\SizeCost$ is replaced by
$\TrainCost$.

Here $\widetilde{\MRE}$ and the tilded cost term are each min-max scaled to
$[0,1]$ within a context, so error and cost are comparable, and $\lambda$ trades
one against the other: small $\lambda$ is accuracy-heavy, large $\lambda$
cost-heavy. We explore a range of $\lambda$, since no single accuracy-versus-cost
weight fits every setting.

This formulation clarifies how the conventional fixed-$n$ protocol fits our
analysis: when strategies are compared at the same sample size with sample count
as the cost axis, the cost term is constant and the ranking collapses to
MRE: the standard pure-accuracy reading, valid for that question. The extra
assumption appears when that fixed-$n$ comparison is read as a fixed-\emph{budget}
one. A resulting ``flip'' is then either a genuine trade-off (the count winner is
more accurate but more expensive under measured cost) or a \emph{domination} (a
different combination is no less accurate and no more expensive). Pareto fronts
expose the second case without choosing a $\lambda$: when the measured-cost front
contains combinations missing from the count-cost front, a size-only analysis has
hidden choices a practitioner might legitimately prefer.

We count only training acquisition. Validation measurements are part of the
experimental protocol and are the same across strategies. A sampling strategy
controls which training configurations is measured.
